\documentclass[aspectratio=169]{beamer}
\usepackage[utf8]{inputenc}
\usepackage[T2A]{fontenc}
\usepackage[russian]{babel}
\usepackage{graphicx}
\usepackage{booktabs}
\usepackage{algorithm}
\usepackage{algpseudocode}
\usepackage{listings}
\usepackage{xcolor}
\usepackage{csquotes}

\usepackage[backend=biber, style=gost-numeric, sorting=none]{biblatex}
\addbibresource{references.bib}

\usetheme{Madrid}
\usecolortheme{default}

\lstdefinestyle{java}{
  language=Java,
  basicstyle=\tiny\ttfamily,
  keywordstyle=\color{blue}\bfseries,
  stringstyle=\color{red},
  commentstyle=\color{gray}\itshape,
  breaklines=true,
  tabsize=2,
  showstringspaces=false
}
\lstset{style=java}

\algrenewcommand\algorithmicif{\textbf{если}}
\algrenewcommand\algorithmicthen{\textbf{тогда}}
\algrenewcommand\algorithmicelse{\textbf{иначе}}
\algrenewcommand\algorithmicend{\textbf{конец}}
\algrenewcommand\algorithmicreturn{\textbf{вернуть}}
\algrenewcommand\algorithmicfunction{\textbf{Функция}}
\algrenewcommand\algorithmicrequire{\textbf{Вход:}}
\algrenewcommand\algorithmicensure{\textbf{Выход:}}

\floatname{algorithm}{Алгоритм}

\title{Тестирование идемпотентности и повторяемости запросов в сетевых API\\
\small по дисциплине <<Методы тестирования программного обеспечения>>}
\author{Плахотников В.\,А.}
\institute{СПбПУ, ИКНТ ВШ ПИ}
\date{2026}

\begin{document}

% ===== СЛАЙД 1: ТИТУЛЬНЫЙ =====
\begin{frame}[plain]
\maketitle
\small
\begin{tabular}[t]{@{}l@{\hspace{1pt}}p{.56\textwidth}@{}}
\qquad \qquad \qquad \qquad \qquad \qquad Выполнил студент группы 5130903/30303:\\
\qquad \qquad \qquad \qquad \qquad \qquad Плахотников В.\,А.
\end{tabular}%

\begin{tabular}[t]{@{}l@{\hspace{1pt}}p{.32\textwidth}@{}}
\\
\qquad \qquad \qquad \qquad \qquad \qquad Руководитель:\\
\qquad \qquad \qquad \qquad \qquad \qquad Старший преподаватель В.\,А.~Пархоменко
\end{tabular}%
\end{frame}

% ===== СЛАЙД 2: СТРУКТУРА =====
\begin{frame}{Структура доклада}
\tableofcontents
\end{frame}

% ===== СЛАЙД 3: ПОСТАНОВКА ПРОБЛЕМЫ =====
\section{Постановка проблемы}
\begin{frame}{Постановка проблемы}
\textbf{Проблема:} При сетевых сбоях клиент повторяет HTTP-запрос (retry).

\vspace{3mm}
\textbf{Без идемпотентности:}
\begin{itemize}
    \item POST-запрос на создание заказа $\times 3$ retry $\Rightarrow$ 3 заказа
    \item POST-запрос на платёж $\times 3$ retry $\Rightarrow$ \textbf{тройное списание}
    \item POST товара с тем же SKU $\Rightarrow$ дубликат в каталоге
\end{itemize}

\vspace{3mm}
\textbf{Цель:} Исследовать и продемонстрировать тестирование идемпотентности сетевых API на примере шести различных типов контроллеров Spring Boot (REST, SOAP, gRPC, GraphQL).
\end{frame}

% ===== СЛАЙД 4: МАТЕМАТИЧЕСКАЯ ФОРМАЛИЗАЦИЯ =====
\section{Математическая формализация}
\begin{frame}{Математическая формализация}
\textbf{Определение.} Операция $O: S \to S$ \emph{идемпотентна}, если:
$$O^n(s) = O(s), \quad \forall n \geq 1, \forall s \in S$$

\vspace{3mm}
\textbf{HTTP-методы (RFC 9110):}

\begin{table}
\centering
\begin{tabular}{@{}lccc@{}}
\toprule
Метод & Safe & Идемпотентный \\
\midrule
GET, HEAD & \checkmark & \checkmark \\
PUT, DELETE & --- & \checkmark \\
POST, PATCH & --- & --- \\
\bottomrule
\end{tabular}
\end{table}

\vspace{3mm}
\textbf{Решение для POST:} Idempotency-Key $k$:
$$O'(s, r, k) = \begin{cases} O(s, r), & k \notin KeyStore \\ cached(k), & k \in KeyStore \end{cases}$$
\end{frame}

% ===== СЛАЙД 5: ПСЕВДОКОД =====
\section{Алгоритм}
\begin{frame}{Алгоритм проверки идемпотентности}
\begin{algorithm}[H]
\footnotesize
\caption{CheckIdempotency($request$)}
\begin{algorithmic}[1]
\If{$request.method \neq$ POST}
    \State \Return ProcessRequest($request$)
\EndIf
\State $key \gets$ headers[``Idempotency-Key'']
\If{$key = \varnothing$}
    \State \Return 400 Bad Request
\EndIf
\State $cached \gets KeyStore.find(key)$
\If{$cached \neq \varnothing$ и $cached.response \neq \varnothing$}
    \State \Return $cached.response$ \Comment{кэшированный ответ}
\EndIf
\State $KeyStore.register(key)$
\State $response \gets ProcessRequest(request)$
\State $KeyStore.save(key, response)$
\State \Return $response$
\end{algorithmic}
\end{algorithm}
\end{frame}

% ===== СЛАЙД 6: ШЕСТЬ ТИПОВ КОНТРОЛЛЕРОВ =====
\section{Шесть типов контроллеров}
\begin{frame}{Шесть кардинально разных типов контроллеров}
\scriptsize
\begin{table}
\centering
\begin{tabular}{@{}p{2.8cm}p{3.3cm}p{2.3cm}p{3.4cm}@{}}
\toprule
\textbf{Тип} & \textbf{Транспорт} & \textbf{Домен} & \textbf{Идемпотентность} \\
\midrule
@RestController        & HTTP/REST (JSON)    & Заказы        & Idempotency-Key (HTTP-фильтр) \\
Functional Endpoints   & HTTP/REST (JSON)    & Платежи       & \emph{тот же} фильтр \\
Spring Data REST       & HTTP/HAL+JSON       & Товары        & @Version + ETag \\
Spring-WS @Endpoint    & SOAP / HTTP+XML     & Отправления   & \emph{тот же} фильтр \\
GraphQL @Controller    & GraphQL / HTTP      & Инвентарь     & \emph{тот же} фильтр \\
@GrpcService           & gRPC / HTTP/2       & Уведомления   & ServerInterceptor \\
\bottomrule
\end{tabular}
\end{table}

\vspace{2mm}
\textbf{Ключевая идея:} четыре из шести транспортов используют \emph{один и тот же} \texttt{IdempotencyFilter}; gRPC --- единственный, требующий отдельного перехватчика, но он делегирует тому же \texttt{IdempotencyService}.
\end{frame}

% ===== СЛАЙД 7: ИЛЛЮСТРАТИВНЫЙ ПРИМЕР =====
\section{Иллюстративный пример}
\begin{frame}{Иллюстративный пример: retry платежа}
\textbf{Сценарий:} Клиент отправляет POST /api/payments, получает таймаут, повторяет запрос.

\vspace{3mm}
\textbf{С идемпотентностью (Idempotency-Key: ``pay-001''):}
\begin{enumerate}
    \footnotesize
    \item Запрос \#1: key=``pay-001'' $\notin$ KeyStore $\Rightarrow$ register, создать платёж, сохранить ответ
    \item Запрос \#2 (retry): key=``pay-001'' $\in$ KeyStore $\Rightarrow$ \textbf{вернуть кэш}
    \item Результат: \textbf{1 платёж} (корректно)
\end{enumerate}

\vspace{3mm}
\textbf{Без идемпотентности:}
\begin{enumerate}
    \footnotesize
    \item Запрос \#1: создать платёж $\Rightarrow$ списание \#1
    \item Запрос \#2 (retry): создать платёж $\Rightarrow$ списание \#2
    \item Результат: \textbf{2 платежа --- двойное списание!}
\end{enumerate}
\end{frame}

% ===== СЛАЙД 8: ПОШАГОВОЕ ВЫПОЛНЕНИЕ =====
\begin{frame}{Пошаговое выполнение на данных}
\footnotesize
\textbf{Вход:} POST /api/orders, body: \{description: ``Ноутбук'', amount: 89999.99\}

\vspace{2mm}
\begin{tabular}{@{}clll@{}}
\toprule
\textbf{Шаг} & \textbf{Действие} & \textbf{KeyStore} & \textbf{orders (count)} \\
\midrule
0 & Начальное состояние & $\varnothing$ & 0 \\
1 & Получен key=``abc'' & $\varnothing$ & 0 \\
2 & find(``abc'') $= \varnothing$ & $\varnothing$ & 0 \\
3 & register(``abc'') & \{abc: pending\} & 0 \\
4 & INSERT order & \{abc: pending\} & \textbf{1} \\
5 & save(``abc'', 201, body) & \{abc: 201, body\} & 1 \\
\midrule
6 & Retry: key=``abc'' & \{abc: 201, body\} & 1 \\
7 & find(``abc'') $\neq \varnothing$ & \{abc: 201, body\} & 1 \\
8 & return cached 201 & \{abc: 201, body\} & \textbf{1} \\
\bottomrule
\end{tabular}

\vspace{3mm}
\textbf{Результат:} 1 заказ, 2-й запрос вернул кэш. Идемпотентность обеспечена.
\end{frame}

% ===== СЛАЙД 9: ТЕСТИРОВАНИЕ =====
\section{Апробация}
\begin{frame}{Результаты тестирования}
\textbf{Технологии:} JUnit 5 + Testcontainers (PostgreSQL в Docker) + MockMvc

\vspace{3mm}
\begin{table}
\centering
\scriptsize
\begin{tabular}{@{}lccc@{}}
\toprule
\textbf{Тестовый класс} & \textbf{Тестов} & \textbf{Транспорт} & \textbf{Результат} \\
\midrule
OrderIdempotencyTest       & 6 & @RestController             & \textcolor{green!70!black}{PASS} \\
PaymentIdempotencyTest     & 5 & Functional                  & \textcolor{green!70!black}{PASS} \\
ProductIdempotencyTest     & 5 & Spring Data REST            & \textcolor{green!70!black}{PASS} \\
SoapIdempotencyTest        & 3 & SOAP (Spring-WS)            & \textcolor{green!70!black}{PASS} \\
GrpcIdempotencyTest        & 4 & gRPC (HTTP/2 + protobuf)    & \textcolor{green!70!black}{PASS} \\
GraphqlIdempotencyTest     & 4 & GraphQL                     & \textcolor{green!70!black}{PASS} \\
DataImportTest             & 2 & JSON/CSV                    & \textcolor{green!70!black}{PASS} \\
\midrule
NonIdempotentProfileTest   & 7 & все 6 транспортов           & \textcolor{green!70!black}{PASS*} \\
\midrule
\textbf{Итого}             & \textbf{36} & & \textbf{PASS} \\
\bottomrule
\end{tabular}
\end{table}

\vspace{2mm}
\footnotesize *NonIdempotentProfileTest: 7 тестов (по одному на каждый транспорт + ранее существующие), все проходят зелёными, \textbf{подтверждая наличие багов} при отключённой идемпотентности.
\end{frame}

% ===== СЛАЙД 10: ДЕМОНСТРАЦИЯ ПРОФИЛЕЙ =====
\begin{frame}{Профиль non-idempotent: демонстрация проблем на 6 транспортах}
\scriptsize
\begin{table}
\centering
\begin{tabular}{@{}p{5cm}p{4.5cm}c@{}}
\toprule
\textbf{Тест} & \textbf{Проблема} & \textbf{Статус} \\
\midrule
REST POST без Idempotency-Key & Проходит без проверки                    & \textcolor{red}{BUG} \\
REST повтор POST с тем же ключом & Дубликат заказа                       & \textcolor{red}{BUG} \\
Data REST дубликат SKU & Проходит (нет constraint)                       & \textcolor{red}{BUG} \\
Functional 3x retry платежа & Тройное списание                           & \textcolor{red}{BUG} \\
SOAP тот же tracking\_number & Два Shipment c одним номером              & \textcolor{red}{BUG} \\
gRPC тот же idempotency-key & Interceptor выкл. $\Rightarrow$ дубликат   & \textcolor{red}{BUG} \\
GraphQL мутация с тем же ключом & Два Inventory с одним item\_name        & \textcolor{red}{BUG} \\
\bottomrule
\end{tabular}
\end{table}

\vspace{2mm}
\textbf{Ключевой вывод:} отключение механизмов идемпотентности приводит к критическим багам \emph{во всех транспортах}, и каждый случай обнаруживается отдельным тестом.
\end{frame}

% ===== СЛАЙД 10а: АРХИТЕКТУРА ПЕРЕИСПОЛЬЗОВАНИЯ =====
\begin{frame}{Архитектура переиспользования}
\scriptsize
Единый \texttt{IdempotencyService} обслуживает шесть транспортов через два механизма перехвата.

\vspace{2mm}
\begin{center}
\begin{tabular}{@{}lll@{}}
\toprule
\textbf{Перехват} & \textbf{Транспорты} & \textbf{Кодировка снэпшота} \\
\midrule
\texttt{IdempotencyFilter} (servlet) & REST, Functional, SOAP, GraphQL & TEXT (JSON / XML) \\
\texttt{GrpcIdempotencyInterceptor}   & gRPC                            & BYTEA (protobuf) \\
\texttt{@Version} + ETag              & Spring Data REST                & --- \\
\bottomrule
\end{tabular}
\end{center}

\vspace{3mm}
\textbf{Доработка фильтра:} миграция V4 расширила \texttt{idempotency\_keys} тремя колонками --- \texttt{response\_content\_type}, \texttt{response\_kind}, \texttt{response\_body\_bytes}. Это превратило фильтр из <<JSON-only>> в \emph{транспорт-агностичный}: SOAP реплеит \texttt{text/xml}, gRPC --- бинарные protobuf-байты через \texttt{MethodDescriptor.responseMarshaller}.
\end{frame}

% ===== СЛАЙД 11: ВЫВОДЫ =====
\section{Выводы}
\begin{frame}{Выводы}
\begin{enumerate}
    \item Формализована математическая модель идемпотентности HTTP: $O^n(s) = O(s)$
    \item Реализованы \textbf{6 типов контроллеров} в 4 транспортах с единым механизмом идемпотентности:
    \begin{itemize}
        \scriptsize
        \item @RestController, Functional, Spring Data REST (HTTP/REST)
        \item Spring-WS @Endpoint (SOAP)
        \item @GrpcService + ServerInterceptor (gRPC HTTP/2)
        \item GraphQL @Controller (single endpoint)
    \end{itemize}
    \item Доработан \texttt{IdempotencyFilter} до транспорт-агностичного: снэпшот ответа $\langle \text{status}, \text{contentType}, \text{kind}, \text{body} \rangle$ позволяет реплеить JSON, XML и бинарный protobuf через единый кеш.
    \item \textbf{36 интеграционных тестов} (Testcontainers + PostgreSQL)
    \item Профиль non-idempotent демонстрирует баги во всех 6 транспортах
\end{enumerate}

\vspace{3mm}
\centering
\textbf{Спасибо за внимание!}
\end{frame}

\end{document}
